\section{The exact algorithm}
\label{sec:algorithm}

An input consists of finite preperiods and cyclic tails for both $(a_i)$ and
$(b_i)$.  This is a finite presentation, but two points require care: legality
must be checked across the cyclic seam, and a plateau of asymptotic length zero
must not be classified by floating-point sampling.

\begin{proposition}[Terminating exact construction]
\label{prop:exact-algorithm}
From a finite ultimately periodic presentation satisfying the hypotheses of
\cref{thm:main}, one can terminate with
\begin{enumerate}[label=\textup{(\roman*)}]
\item the exact finite set $\mathcal C$ in \eqref{eq:finite-maximum};
\item a primitive minimal polynomial and rational isolating interval for
      $\dio(x)$; and
\item an exact equality and order decision for the outputs of two inputs.
\end{enumerate}
\end{proposition}

\begin{proof}
We give the construction in steps.

\smallskip
\noindent\emph{Input validation.}
Check $a_i\geq1$, the exceptional first bound $0\leq b_1\leq a_1-1$, and
all later bounds and carry implications in \eqref{eq:legality}.  Because the
conditions are local, the finite preperiod followed by two copies of the tail
also checks the periodic seam.  Enlarge and align the two tails to a common
period $L$.  The case of each phase depends only on the adjacent displayed
digits, so the same finite check selects its rows in \cref{tab:eight-cases}.

\smallskip
\noindent\emph{Exact spectral data.}
Construct the integer matrices $M$ and $P$ in
\eqref{eq:block-period}.  The Perron root is the larger real root of
\begin{equation}
 f(X)=X^2-\tr(P)X+(-1)^L.
 \label{eq:perron-polynomial}
\end{equation}
Factor $f$ over $\Q$ and use Sturm-sequence root isolation to select its larger
real root.  This also covers a rational root by the degree-one branch.  Thus
$K=\Q(\Lam)$ is represented exactly, with each element reduced to
$x+y\Lam$, $x,y\in\Q$.  The projectors in the proof of
\cref{lem:three-scales} compute the three coefficients of every endpoint,
denominator, and interval length.

\smallskip
\noindent\emph{Eventual nonemptiness.}
For a nominal plateau let $D_k=B_k-A_k$.  It is nonempty exactly when
$D_k\geq0$.  Write its exact phase expansion as
\begin{equation}
 D_k=C_\Lam\Lam^k+C_1+C_\mu\mu^k,
 \qquad |\mu|<1.
 \label{eq:length-expansion}
\end{equation}
If $C_\Lam\ne0$, its exact sign decides the eventual sign of $D_k$.  If
$C_\Lam=0$ and $C_1\ne0$, the exact sign of $C_1$ does so.  If both vanish,
$D_k$ is an integer sequence tending to zero; hence $D_k=0$ for all
sufficiently large $k$, and the row is an eventual singleton rather than an
empty interval.  This also handles odd $L$, when $\mu<0$: an alternating
nonzero tail cannot remain an integer while tending to zero.

Zero and sign in $K$ are decidable.  One may clear rational denominators and
use the isolating interval for $\Lam$; exact zero is tested algebraically
before interval refinement, and refinement terminates for a nonzero element.
This produces the sets $\mathcal R_j$ and hence the at most $4L$ elements of
$\mathcal C$.

\smallskip
\noindent\emph{Maximum and certified output.}
Compare two candidates by the exact sign of their difference in $K$, retaining
one representative when the difference is zero.  The winning value and hence
$\dio(x)$ reduce to $z=x+y\Lam$ with $x,y\in\Q$.  If $y=0$, its minimal
polynomial is linear.  Otherwise conjugating the quadratic equation
\eqref{eq:perron-polynomial} gives a rational quadratic polynomial for $z$;
clearing denominators, taking the irreducible factor containing $z$, and
normalizing its content and leading sign gives the primitive minimal
polynomial.  Rational interval arithmetic, refined against the isolated root
$\Lam$, produces a rational interval containing $z$ and no other root.

For outputs from two possibly different quadratic fields, compute the gcd of
their primitive minimal polynomials.  If it is $1$, the numbers are unequal.
If a common factor remains, compare their isolating intervals against the
isolated roots of that factor; this decides equality.  Only after inequality
is known do we refine the two intervals until they are disjoint, which decides
order.  Therefore an equal pair cannot cause an endless numerical refinement.
Every operation above is an exact terminating operation over $\Z$ and $\Q$.
\end{proof}

\begin{proof}[Proof of \cref{thm:main}]
The complete corrected table and \eqref{eq:dio-repetition}, through
\cref{prop:finite-interface}, reduce the exponent to plateau endpoint ratios.
\Cref{lem:three-scales,lem:four-per-phase} and
\eqref{eq:finite-maximum} put at most $4L$ of their limits in the single field
$\Q(\Lam)$.  A maximum of finitely many elements is one of them, so adding
$1$ preserves algebraic degree at most two.  The construction, output, and
comparison assertions are exactly \cref{prop:exact-algorithm}.
\end{proof}
